% Unregistered candidate geometry.  No Chow or negative-K conclusion.
\paragraph{An explicit ordered small resolution of a tetrahedral fibre.}
Let $k$ be algebraically closed of characteristic $p>0$.  In four
variables we put
\begin{equation}
 Q=x_0x_1x_2x_3,\qquad
 G=\begin{cases}
 \sum_{i=0}^3x_i^5,&p\ne5,\\
 \sum_{i=0}^3x_i^5+\sum_{0\le i<j\le3}x_i^4x_j,&p=5,
 \end{cases}
 \qquad F=Q-G.
 \tag{TG.1}
\end{equation}
Let $R$ be the local ring of $F=0$ at the origin, or its strict
henselization.  We construct a projective morphism $Z\to\operatorname{Spec}R$
which is an isomorphism off the closed point and whose source is regular.
Its reduced closed fibre $E$ is an effective Cartier simple normal
crossings divisor with four components.  In the order $0,1,2,3$, these
components are respectively blowups of $\mathbf P^2_k$ at $15,10,5,0$
distinct points on their coordinate edges.  Every double curve is
$\mathbf P^1_k$, the four triple intersections are single points, and
there is no quadruple intersection.  Thus its dual complex is the
boundary of a tetrahedron.

\emph{Step 1: the isolated local hypersurface and the first blowup.}
Every coordinate-plane section of $G=0$ is smooth, every coordinate
edge meets it in five distinct nonvertex points, and $G$ is nonzero at
the four coordinate vertices.  For $p\ne5$ these assertions follow
directly from the Fermat equation.  For $p=5$, the plane polynomial is
\[
 g=a^5+b^5+c^5+a^4b+a^4c+b^4c.
\]
At a common zero of $g$ and its three derivatives, $a=0$ would imply
$b=0$ and then $g=c^5\ne0$.  If $a\ne0$, the first derivative gives
$c=-b$, and the third gives $a^4+b^4=0$.  Substitution in $g=0$ gives
$a^5=b^5$, hence $a=b$.  This contradicts $2a^4=0$ in characteristic
$5$.  On an edge the polynomial is $a^5+b^5+a^4b$; in the chart $b=1$
its derivative is $4a^3$ and its value at $a=0$ is $1$.  There are
therefore five distinct roots, none at a coordinate vertex.

Euler's homogeneous identities give
\[
 \sum_i x_iF_{x_i}=4Q-5G.
\]
At a singular point of $F=0$ we consequently have $Q=G=0$, in every
characteristic.  A nonzero such point lies on a coordinate face.  The
restriction of $Q$ to that face is identically zero, so the tangential
derivatives of $G$ there would all vanish.  The smooth plane and edge
sections just checked exclude this, and the vertices have $G\ne0$.
Thus the affine hypersurface is smooth away from the origin.  At the
origin its local ring is a domain: a factorization into two nonunits
in the ambient regular local UFD would make their intersection, of
dimension at least two, singular.  Repeated factors are excluded for
the same reason.  The hypersurface is Cohen--Macaulay and regular in
codimension one, hence normal.

The blowup $Z_0$ of the origin is the strict transform inside
$\operatorname{Tot}(\mathcal O_{\mathbf P^3}(-1))$.  With $t$ the fibre
coordinate and $w_i$ homogeneous base coordinates, its equation is
\begin{equation}
 w_0w_1w_2w_3=tG(w_0,w_1,w_2,w_3).
 \tag{TG.2}
\end{equation}
This description is the actual Rees blowup, obtained on each chart by
cancelling the common fourth power in the domain.  Its Cartier closed
fibre $E_0=(t=0)$ is the reduced union of the four coordinate planes.
At a triple vertex $G$ is a unit, so $t$ is a unit times the product of
three regular parameters.  Away from the vertices and double edges
the source is regular.  On a double edge it is regular unless $G=0$.
At each of the thirty zeros of $G$ on the six edges, its simple zero
permits an etale change of coordinates putting (TG.2) into the form
\begin{equation}
 uv=tw,\qquad E_0=(t=0),\qquad
 (E_i,E_j)=((u,t),(v,t)).
 \tag{TG.3}
\end{equation}
Indeed, after dividing by the two nonzero homogeneous coordinates,
we may use the edge parameter $G$ as one of the local coordinates.
There are no other singularities of $Z_0$.

\emph{Step 2: the two actual blowup charts.}
For the domain in (TG.3), blow up the ideal $(u,t)$.
The two affine Rees charts are
\begin{equation}
 \begin{array}{c|c|c}
 \text{chart}&\text{independent coordinates}&\text{map to (TG.3)}\\ \hline
 u&u,a,w&t=ua,\quad v=aw,\\
 t&t,b,v&u=tb,\quad w=bv.
 \end{array}
 \tag{TG.4}
\end{equation}
Here $a=t/u$ and $b=u/t$, with $ab=1$ on the overlap.
The cancellation giving $v=aw$ or $w=bv$ takes place in the affine
blowup algebra inside the fraction field; we do not retain the extra
component of an unsaturated total-transform equation.
Both charts are regular, and the exceptional fibre over the node is
one $\mathbf P^1$.  In the first chart the inverse image of $E_0$ is
$(ua=0)$; in the second it is $(t=0)$.  It is therefore reduced with
simple normal crossings.  The chosen plane $(u,t)$ becomes its
ordinary point blowup, with charts $v=aw$ and $w=bv$; the other plane
is unchanged.  Their double curve is its strict transform and maps
isomorphically to the original edge.  In particular the operation is
small: it introduces no new divisorial component.

\emph{Step 3: global projective compatibility.}
We order the four planes and successively blow up the ideal sheaves of
their strict transforms.  Each is a global closed reduced Weil
divisor.  On the regular locus it is locally Cartier, so its blowup
is the identity.  At an unresolved node on an incident edge the
operation is exactly (TG.4); at a node on another edge it is the
identity.  Once one endpoint plane of an edge has been processed,
all five nodes on that edge have been resolved.  Every later centre
is Cartier near them and creates no further change.  The triple
vertices are regular throughout and are unchanged.

These facts use the actual affine blowup algebras, flat base change,
and the identity for a Cartier centre in
\citet[Lemmas 31.33.2--31.33.4 and 31.33.7, Tags 0804, 0805, 02OS,
and 0807]{rootStacksBlowup}.  Every ideal is coherent, so all the
blowups are projective by Lemma 31.33.13, Tag 02NS, in the same source.
After processing the first three planes no nodes remain; the fourth
blowup is the identity.  The resulting $Z$ is regular and projective
over $R$, and it preserves the puncture.  The node sets assigned to
the earlier endpoint have respective sizes $15,10,5,0$.  All points
are off the coordinate vertices.  This proves the asserted component
and incidence description.  Localization and strict henselization
are flat; the same charts and construction therefore apply to both
versions of $R$.

\emph{Step 4: the two distinct normal-positivity statements.}
Write $\rho:E\to E_0\subset\mathbf P^3$ for the induced morphism.
The full closed fibre remains reduced and Cartier under (TG.4), so
\begin{equation}
 \mathfrak m_R\mathcal O_Z=\mathcal O_Z(-E),\qquad
 \mathcal O_Z(-E)|_E=\rho^*\mathcal O_{E_0}(1).
 \tag{TG.5}
\end{equation}
In particular, on each plane blowup the conormal is the pullback of
$\mathcal O_{\mathbf P^2}(1)$.  It is nef and has degree zero on each
of the thirty small exceptional curves; it is not ample.

There nevertheless exists a positive integral combination
$D=\sum_i d_iE_i$ with $-D$ ample over $\operatorname{Spec}R$.
Initially $\mathcal O_{Z_0}(-E_0)$ is relatively ample.  At a nontrivial
ordered blowup $h$, the local chart calculation gives
\[
 \mathcal I_{E_i}\mathcal O_{Z'}=\mathcal O_{Z'}(-E_i'),
\]
which is its relatively ample tautological sheaf.  If $-D_{\mathrm{old}}$
is ample over the original base, then
$\mathcal O(-E_i')\otimes h^*\mathcal O(-D_{\mathrm{old}})^{\otimes M}$
is ample over that base for $M\gg0$, by
\citet[Lemma 29.38.8, Tag 0C4K]{frC3beStacksCompositeAmple}.
Its negative divisor is $E_i'+M h^*D_{\mathrm{old}}$.
The morphism is small, so pulling back the already Cartier divisor
$D_{\mathrm{old}}$ retains its positive coefficients on the four
strict components.  Induction constructs the required $D$.
This weighted ample divisor is not substituted for the fixed
conormal in (TG.5).

\paragraph{Scope of the construction.}
We have verified the regular projective geometry, the strict-local
base change, the reduced SNC closed fibre, its dual complex, and its
conormal.  We have not shown that a dual-complex $K_2(k)$ coefficient
lifts to $K_0(Z)$, that a proper-curve relation lifts to $Z\times P$,
or that an ordinary top Chow quotient is nonzero.  Those questions
retain the exact negative-$K$ boundary described in (FF.1)--(FF.3).
